\section{Independent verification of residue-section
selection}\label{independent-verification-of-residue-section-selection}

25 September 2026. Complete review and further derivation RSR0--RSR8.

The reviewed source is
\nolinkurl{RESIDUE_SECTION_SELECTION_AND_ORIGINAL_IDEAL.md}, RSS0--RSS6,
SHA256
\nolinkurl{ac69b79cac80878918e8a482f26d2bed7d7d69fd8ce3267ba3bee751f5ac61a7}.
This review verifies the complete source, not a summary. No mathematical
formula error was found. The single wording correction reported to the
author was to specify the Jensen center as a point \(w_0\) with
\(L(w_0)\ne0\). The reviewed edition includes that correction.

\subsection{RSR0. Original objects, construction stage and
coverage}\label{rsr0.-original-objects-construction-stage-and-coverage}

The supporting datum remains
\(\tau\langle Z_1;\mathrm{no}\ Z_2\rangle\). Its addition remains
retracted. No parity, coordinate, arithmetic value, metric or
coefficient operation is assigned to it here. The scalar arithmetic,
recovered positive integers, logarithms and analytic variable below are
used after the complete-history arithmetic reconstruction. Separate
branch histories remain separate.

Retain the original complete Fréchet source \[
\mathcal B=\{F\in\mathcal O(\mathbb C):q_{A,N}(F)<\infty
\text{ for every }A>0,N\in\mathbb Z_{\ge0}\},\qquad
q_{A,N}(F)=\sup_{|\Re s|\le A}(1+|\Im s|)^N|F(s)|.
\tag{RSR0.1}
\] The original ideal \(I\) consists of every \(F\in\mathcal B\)
satisfying

\[
F^{(j)}(\rho)=0\quad(0\le j<m_\rho)
\tag{RSR0.2}
\]

at every actual nontrivial zero \(\rho\) of original \(\zeta\), with its
full multiplicity \(m_\rho\). These jet maps are continuous: Cauchy's
derivative estimate on a fixed circle about each point bounds them by a
source seminorm on a larger strip. Consequently \(I\) is closed. Write
\(Q=\mathcal B/I\).

The complete original comparison and all endpoint data remain

\[
F_0(s)=\frac{s(s-1)}8\pi^{-s/2}\Gamma(s/2)\zeta(s),\qquad
F_0(0)=F_0(1)=\frac18,\qquad F_0(-1)=F_0(2)=\frac\pi{24},
\] \[
F_0(-2j)=\frac{j(2j+1)(-1)^j\pi^j}{2j!}\zeta'(-2j)
\quad(j\ge1).
\tag{RSR0.3}
\]

The notation \(F_0\) is the full source multiplier, not a replacement
for original zeta. Its Gaussian comparison is \(g_t(s)=e^{ts^2}\), with
the same fixed real \(t>0\) throughout. The actual meromorphic functions
are

\[
h_t=\frac{8g_tF_0}{s},\qquad j_t=\frac{g_t}{s},\qquad
E=\mathcal B+\mathbb C h_t.
\tag{RSR0.4}
\]

The topology of \(E\) is transported from its unique coordinates
\(F+ch_t\), with \(c=\operatorname{res}_0(F+ch_t)\). The cover action is
\(U_nf(s)=n^{1-s}f(s)\), for all recovered positive integers. Its
multiplier and reciprocal are bounded on every fixed strip, so its
restriction is a continuous automorphism of \(\mathcal B\).

This review read RSS0--6, NPE0--10, NER0--11 and SSR0--7 completely. The
full-source injectivity proof NHJ0--9 and original ideal proof NCI0--8
were previously completely derived and checked in this same independent
mathematical task; their exact receiving statements are identified
below. Their reading and hash receipt accompanies this review. The
original Hardy coefficients come from S. Waleed Noor,
\href{https://arxiv.org/abs/1809.09577v4}{arXiv:1809.09577v4},
original-author TeX SHA256
\nolinkurl{bc3075483547782dce36bf55a6349899a26766fcfc8ae9f265079ec74601f2d3},
canonical source ID \nolinkurl{PUBUNIT-803463A3EF787AE3E69C519B}; the
earlier source reading in this task covered lines 1--387, including the
relevant original proofs. No new entire-paper reading is claimed. The
comparison with \href{https://arxiv.org/abs/0903.2024v3}{Connes--Consani
§5} and \href{https://numdam.org/item/PMIHES_1980__52__137_0/}{Deligne,
Weil II §3.6} remains a programme target, not an assumed identification
used in this proof.

\subsection{RSR1. The logarithmic sampling proof works in the original
topology}\label{rsr1.-the-logarithmic-sampling-proof-works-in-the-original-topology}

Define the actual entire discrepancy

\[
\eta_{n,t}=nj_t-U_nj_t=g_t(s)\frac{n-n^{1-s}}s,
\qquad \eta_{n,t}(0)=n\log n.
\tag{RSR1.1}
\]

For \(w\in\mathbb C\), retain both equal expressions

\[
K_w(s)=g_t(s)\frac{e^w-e^{(1-s)w}}s
=we^wg_t(s)\int_0^1e^{-\theta sw}\,d\theta,
\qquad K_w(0)=we^w.
\tag{RSR1.2}
\]

The integral identity proves removability at \(s=0\). Put \(s=x+iy\) and
\(w=a+ib\). For \(0\le\theta\le1\), the modulus of the exponential in
the integrand, including \(e^wg_t\), is bounded on \(|x|\le A\) by

\[
e^{tA^2+(1+A)|a|}e^{-ty^2+\theta by}
\le e^{tA^2+(1+A)|a|+b^2/(4t)}
e^{-t(y-\theta b/(2t))^2}.
\tag{RSR1.3}
\]

Writing \(y=(y-\theta b/(2t))+\theta b/(2t)\) bounds the weighted
supremum by

\[
q_{A,N}(K_w)\le C_{A,N,t}|w|(1+|b|)^N
e^{(1+A)|a|+b^2/(4t)}.
\tag{RSR1.4}
\]

Each fixed \(w\)-derivative introduces finitely many factors of \(s\),
\(w\), and bounded powers of \(\theta\). On compact parameter sets the
same negative vertical quadratic controls these factors in every source
seminorm. The integral Taylor remainder obeys these estimates too; hence
\(w\mapsto K_w\) is entire as a \(\mathcal B\)-valued map, not merely
pointwise entire.

Let \(\Lambda\in\mathcal B'\) annihilate every \(\eta_{n,t}\),
\(n\ge2\). A continuous functional is bounded by finitely many defining
seminorms; enlarging \(A,N\) combines them. Thus

\[
L(w)=\Lambda(K_w),\qquad |L(w)|\le C e^{C(1+|w|^2)}.
\tag{RSR1.5}
\]

Here \(L(\log n)=0\) for every \(n\ge2\), and \(L(0)=0\). If \(L\) were
nonzero, choose \(w_0\) with \(L(w_0)\ne0\). Jensen's formula on a disk
centered at \(w_0\), with outer radius \(2R\), bounds the number of
zeros in radius \(R\), counted with multiplicity, by \(C_0+C_1R^2\):
each such zero contributes at least \(\log2\), whereas the circle
maximum is bounded by (RSR1.5). Radii meeting boundary zeros may be
replaced by radii between \(2R\) and \(3R\), with the same conclusion.
For sufficiently large \(R\), every distinct point \(\log n\),
\(2\le n\le e^{R/2}\), lies in \(|w-w_0|<R\). Their count is exponential
in \(R\), contradicting the quadratic bound. Therefore \(L=0\).

The differentiated identity, including the sign, is

\[
(\partial_w-1)K_w=g_t(s)e^{(1-s)w}.
\tag{RSR1.6}
\]

Indeed the derivative and subtraction leave \(s e^{(1-s)w}\) in the
numerator. The same identity holds at \(s=0\) by removability. Setting
\(w=-v\) and dividing the nonzero scalar \(e^{-v}\) yields

\[
\Lambda(g_te^{vs})=0\qquad(v\in\mathbb C).
\tag{RSR1.7}
\]

This uses the whole recovered cover family; it does not initialize
arithmetic using selected degrees.

\subsection{RSR2. Exact half-Mellin passage from the tests to the full
source}\label{rsr2.-exact-half-mellin-passage-from-the-tests-to-the-full-source}

Fix \(F\in\mathcal B\) and real \(u>t\). Put
\(H=e^{(u-t)s^2}F\in\mathcal B\). The original inverse transform is

\[
a_H(x)=\frac{x^{-1/2}}\pi\int_{\mathbb R}H(1/2+iy)x^{-iy}\,dy.
\tag{RSR2.1}
\]

Cauchy's derivative estimates on small circles show that
\(y\mapsto H(c+iy)\), together with every derivative, is rapidly
decreasing for every fixed real \(c\). Fourier inversion is therefore
applicable. With \(v=\log x\), the original constant is retained
exactly:

\[
H(s)=\frac12\int_{\mathbb R}a_H(e^v)e^{sv}\,dv.
\tag{RSR2.2}
\]

For clarity, (RSR2.1) is \(1/\pi\) times the inverse Fourier integral of
\(H(1/2+iy)\), followed by the factor \(e^{-v/2}\). The Fourier
inversion constant \(1/(2\pi)\) consequently makes the forward integral
in (RSR2.2) have coefficient \(1/2\).

For each fixed real \(c\), contour shifting in (RSR2.1) gives

\[
a_H(e^v)=\frac{e^{-cv}}\pi
\int_{\mathbb R}H(c+iy)e^{-iyv}\,dy.
\tag{RSR2.3}
\]

The horizontal edges tend to zero because their real coordinates are in
a fixed compact interval and \(H\) has rapid vertical decay there. For
every \(M>0\), take \(c=M\) when \(v\ge0\) and \(c=-M\) when \(v<0\).
The integral absolute values are bounded by the corresponding
rapid-decay seminorms. Hence

\[
|a_H(e^v)|\le C_{H,M}e^{-M|v|}.
\tag{RSR2.4}
\]

This proves absolute convergence of (RSR2.2) at every \(s\), locally
uniformly, so equality on the original vertical line extends to all
\(s\). For real \(v\),

\[
q_{A,N}(g_te^{vs})\le C_{A,N,t}e^{A|v|}.
\tag{RSR2.5}
\]

Thus the following integral converges in every original source seminorm,
by (RSR2.4), and in the complete source itself:

\[
e^{us^2}F(s)=\frac12\int_{\mathbb R}
a_H(e^v)g_t(s)e^{sv}\,dv.
\tag{RSR2.6}
\]

Finite-interval Riemann sums converge in each seminorm and the tails
tend to zero by the displayed bounds. A continuous functional passes
through this integral. Equation (RSR1.7) implies
\(\Lambda(e^{us^2}F)=0\) for every real \(u>t\).

For fixed \(F\), the \(\mathcal B\)-valued function
\(u\mapsto e^{us^2}F\) is holomorphic on \(\Re u>0\). Indeed, writing
\(u=p+iq\) gives

\[
\Re(us^2)=p(x^2-y^2)-2qxy.
\tag{RSR2.7}
\]

On compact subsets of that half-plane, \(p\) is bounded below by a
positive number and \(|q|\) is bounded above. For \(|x|\le A\), the
resulting uniformly negative quadratic in \(y\) controls every parameter
derivative and Taylor remainder in every seminorm. The scalar identity
theorem now gives zero for all \(\Re u>0\). Finally, for real
\(0<u\le1\), the integral remainder

\[
e^{us^2}-1=us^2\int_0^1e^{\theta us^2}\,d\theta
\]

gives exactly

\[
q_{A,N}((e^{us^2}-1)F)
\le u e^{A^2}(A^2+1)q_{A,N+2}(F).
\tag{RSR2.8}
\]

Letting \(u\downarrow0\) proves \(\Lambda(F)=0\). The continuous
annihilator of the discrepancy span is therefore zero. Hahn--Banach
separation of closed subspaces in this Hausdorff locally convex space
proves

\[
\boxed{\overline{\operatorname{span}_{\mathbb C}\{\eta_{n,t}:n\ge2\}}^{\mathcal B}=\mathcal B.}
\tag{RSR2.9}
\]

This completes the RSS1 proof on the full original topology, with no
polynomial-density assumption.

\subsection{RSR3. Polynomial division, topology and every prescribed
jet}\label{rsr3.-polynomial-division-topology-and-every-prescribed-jet}

Let \(P\in\mathbb C[s]\) satisfy \(P(0)=1\). The constant polynomial
\(1\) is included. Let \(a\ne0\) run over its distinct roots, with
orders \(d_a\), and put \(d=\deg P=\sum_a d_a\). The residue-one section
and its discrepancy are exactly

\[
h_{P,t}=\frac{g_tP}{s},\qquad
h_{P,t}-j_t=\frac{g_t(P-1)}s\in\mathcal B,
\qquad \delta^P_{n,t}=P\eta_{n,t}.
\tag{RSR3.1}
\]

The first difference is a Gaussian times a polynomial, since
\(P(0)-1=0\). Thus all these sections belong to the same \(E\), with the
same residue and topology. Polynomial multiplication satisfies
\(q_{A,N}(PF)\le C_{P,A}q_{A,N+d}(F)\).

Write \(J_P\) for the source subspace with \(F^{(j)}(a)=0\) for all
\(j<d_a\). The product rule gives \(P\mathcal B\subset J_P\). For the
reverse inclusion \(F/P\) is entire, since all root orders are retained.
Choose small disjoint closed disks about the finitely many roots. On the
complement of their interiors, \(|1/P|\) is bounded by some \(C_P\): for
large \(|s|\) use the nonzero leading coefficient; on the remaining
compact set use its nonzero minimum. On the boundary of each disk,

\[
|F/P|\le C_P\sup|F|.
\]

The maximum principle for the entire quotient extends this bound inside
the disk. All these disks lie in a fixed compact set, where the vertical
weight is bounded. Enlarging one fixed strip to contain them gives, for
every \(A,N\),

\[
q_{A,N}(F/P)\le C_Pq_{A,N}(F)
+C_{P,A,N}q_{A_0,0}(F).
\tag{RSR3.2}
\]

For \(P=1\) the disk contribution is absent. This proves the exact
identity and a continuous inverse to multiplication:

\[
\boxed{P\mathcal B=J_P,\qquad
m_P:\mathcal B\longrightarrow P\mathcal B
\text{ is a topological isomorphism}.}
\tag{RSR3.3}
\]

The image is closed because \(J_P\) is a finite intersection of
continuous jet kernels. Multiplying approximations in (RSR2.9) by \(P\),
and using that closed image, proves

\[
\boxed{\overline{\operatorname{span}_{\mathbb C}\{\delta^P_{n,t}:n\ge2\}}^{\mathcal B}=P\mathcal B.}
\tag{RSR3.4}
\]

To identify the full quotient, let \(J(F)=(F^{(j)}(a))_{a,j<d_a}\). It
is continuous and onto \(\bigoplus_a\mathbb C^{d_a}\). A detailed
construction is as follows. Multiplication by \(e^{s^2}\) acts on the
finite jet coordinates at each \(a\) by the lower triangular matrix

\[
(p^{(b)}(a))_{b<d_a}\longmapsto
\left(\sum_{b=0}^j\binom jb(e^{s^2})^{(j-b)}(a)p^{(b)}(a)\right)_{j<d_a}.
\tag{RSR3.5}
\]

Its diagonal entry is the nonzero \(e^{a^2}\), so it is invertible.
Prescribe the inverse jets of a polynomial \(p\). The ideals
\((s-a)^{d_a}\) are pairwise coprime; polynomial Bezout identities give
the Chinese remainder map onto their finite Taylor classes, with the
Taylor coefficient \(p^{(j)}(a)/j!\) retained. Thus such a polynomial
exists. Then \(e^{s^2}p\in\mathcal B\) realizes the original jets.
Choosing the polynomial representatives linearly gives a continuous
linear right inverse from this finite-dimensional space to
\(\mathcal B\). The kernel is exactly \(P\mathcal B\). Consequently

\[
\mathcal B/(P\mathcal B)\simeq\bigoplus_a\mathbb C^{d_a}
\tag{RSR3.6}
\]

is a topological isomorphism, not just a vector-space bijection. Its
complete continuous dual is precisely the span of the functionals
\(F\mapsto F^{(j)}(a)\), \(j<d_a\), with no jet removed.

\subsection{RSR4. Transpose signs, maximal invariant slice and full jet
action}\label{rsr4.-transpose-signs-maximal-invariant-slice-and-full-jet-action}

In the product coordinates \(f=F+ch_{P,t}\), the actual cover action is

\[
U_n^E(F,c)=(U_nF-c\delta^P_{n,t},nc).
\tag{RSR4.1}
\]

A continuous complex-linear functional has coordinates

\[
\widetilde\Lambda(F+ch_{P,t})=\Lambda(F)+c\alpha_P.
\]

Direct substitution, without complex conjugation in this transpose,
gives

\[
(U_n^E)'(\Lambda,\alpha_P)
=(U_n'\Lambda,n\alpha_P-\Lambda(\delta^P_{n,t})).
\tag{RSR4.2}
\]

If a vector in the slice \(\alpha_P=0\) has every first cover image in
that slice, it annihilates every \(\delta^P_{n,t}\), and hence
\(P\mathcal B\) by (RSR3.4). Conversely, \(P\mathcal B\) is invariant
under \(U_n\), since the two multipliers commute. Its annihilator is
preserved by \(U_n'\), and the last coordinate of (RSR4.2) vanishes
there. This proves the exact maximality statement, including a stronger
pointwise formulation:

\[
\boxed{\{(\Lambda,0)_P:(U_n^E)'(\Lambda,0)_P
\text{ has last coordinate }0\text{ for every }n\ge2\}
=(P\mathcal B)^\perp\oplus\{0\}.}
\tag{RSR4.3}
\]

The set on the right is itself invariant. The full action on each jet is

\[
U_n'\operatorname{ev}_a^{(j)}
=n^{1-a}\sum_{b=0}^j\binom jb(-\log n)^{j-b}
\operatorname{ev}_a^{(b)}.
\tag{RSR4.4}
\]

This is exactly the Leibniz formula for \((n^{1-s}F)^{(j)}(a)\). It
preserves the original degree \(n\), complex exponent, every
multiplicity and all lower derivative terms. For \(P=1\), the quotient
has dimension zero and (RSR4.3) is zero. For a general \(P\), it is the
complete finite jet dual from (RSR3.6).

The quotient source is equally explicit:

\[
E/(P\mathcal B)\longrightarrow
\mathcal B/(P\mathcal B)\oplus\mathbb C_\chi,
\qquad[F+ch_{P,t}]\longmapsto([F],c),\quad\chi(n)=n.
\tag{RSR4.5}
\]

It is a topological isomorphism by the product quotient, and its cover
action is diagonal because \(\delta^P_{n,t}\in P\mathcal B\). Thus
\(c\mapsto[ch_{P,t}]\) is an equivariant residue section in this
quotient. This statement is not a section in \(E\) itself. An
equivariant section there would require \((U_n-n)F=\delta^P_{n,t}\),
impossible at zero: the two sides have respective values \(0\) and
\(n\log n\).

\subsection{RSR5. The exact shear retains the original source and
receiver}\label{rsr5.-the-exact-shear-retains-the-original-source-and-receiver}

Define the actual entire differences

\[
k=j_t-h_t=\frac{g_t(1-8F_0)}s,\qquad
b_P=h_{P,t}-h_t=\frac{g_t(P-8F_0)}s.
\tag{RSR5.1}
\]

Their numerators belong to \(\mathcal B\) and vanish at zero. For a
source element \(K\) with \(K(0)=0\), division outside the unit disk and
the maximum principle on the radius-two disk give

\[
q_{A,N}(K/s)\le q_{A,N}(K)+2^{N-1}q_{2,0}(K).
\tag{RSR5.2}
\]

Hence both differences are in \(\mathcal B\). If
\(f=F_h+ch_t=F_P+ch_{P,t}\), then

\[
F_P=F_h-cb_P,
\qquad \alpha_P=\alpha_h+\Lambda(b_P).
\tag{RSR5.3}
\]

These are inverse continuous shears with their opposite signs in the
primal and dual coordinates. On the strong dual, evaluation at \(b_P\)
is continuous because its singleton is a bounded source set. No
underlying function or topology changes.

Let \(\delta_{n,t}=8F_0\eta_{n,t}\in I\) be the original discrepancy.
Then

\[
\delta^P_{n,t}=\delta_{n,t}+(n-U_n)b_P.
\tag{RSR5.4}
\]

Consequently the very same original invariant sector is, in the
\(P\)-coordinates,

\[
\{(\Lambda,\Lambda(b_P))_P:\Lambda\in I^\perp\}.
\tag{RSR5.5}
\]

Its invariance follows without a deleted term:

\[
n\Lambda(b_P)-\Lambda(\delta^P_{n,t})
=\Lambda(U_nb_P)=(U_n'\Lambda)(b_P),
\tag{RSR5.6}
\]

since \(\Lambda(\delta_{n,t})=0\). For \(P=1\) this is exactly RSS4's
nonzero-coordinate graph with \(b_P=k\), even though the new
zero-coordinate invariant slice is zero.

Retain every original Noor coefficient,

\[
\phi_0(s)=-1/s,\qquad
\phi_m(s)=\frac{m^{1-s}-(m+1)^{1-s}}s\quad(m\ge1),
\qquad G_m=g_t\phi_m.
\tag{RSR5.7}
\]

Each \(G_m\in E\) has residue \(-1\). Its corrected entire
representatives are \(G_m+h_t\) and \(G_m+h_{P,t}\), differing by
\(b_P\). The source seminorms have polynomial growth in \(m\): for
\(m\ge1\), the exact integral of the derivative of \(x^{1-s}\) gives

\[
\phi_m(s)=-\frac{1-s}{s}\int_m^{m+1}x^{-s}\,dx.
\tag{RSR5.8}
\]

Outside a fixed disk the Gaussian times this expression has a polynomial
\(m\)-bound in every strip seminorm; the corrected representative is
entire, so inside the disk its boundary bound supplies the same
conclusion. The fixed added section difference does not change
polynomial growth. Hence the receiver series converge on every compact
subdisk, continuously on the strong dual, and give

\[
H_P\Lambda=H_h\Lambda+\frac{\overline{\Lambda(b_P)}}{1-z},
\qquad
M_t\widetilde\Lambda
=H_h\Lambda-\frac{\overline{\alpha_h}}{1-z}
=H_P\Lambda-\frac{\overline{\alpha_P}}{1-z}.
\tag{RSR5.9}
\]

The conjugation belongs to the receiver series, not to the
complex-linear transpose in (RSR4.2). Formula (RSR5.9) is equality of
actual functions, so the shear is a homeomorphism of the full Hardy
graphs, with identical Hardy image and norm. It asserts no separate
Hardy membership of terms involving \((1-z)^{-1}\), whose nonzero
constant coefficient sequence is not square summable. NPE6--7/NER7--8
give injectivity, closed graph and covariance of this full receiver;
none of these properties changes under (RSR5.3).

The retained degree equation is consequently still

\[
n\|M_t\widetilde\Lambda\|^2-
\|M_t(U_n^E)'\widetilde\Lambda\|^2
=n\|(I-P_n)M_t\widetilde\Lambda\|^2.
\tag{RSR5.10}
\]

Here \(W_n\) repeats each coefficient \(n\) times, \(W_n^*W_n=nI\),
\(W_nW_n^*=nP_n\), and \(P_n\) is block averaging. Substituting the
exact covariance into the norm difference proves (RSR5.10). The source
shear removes no degree or projection term.

\subsection{RSR6. Original full ideal and the precise quotient
topology}\label{rsr6.-original-full-ideal-and-the-precise-quotient-topology}

The polynomial argument has not been substituted for the original
nonpolynomial one. The accepted NCI2--5 theorem on this exact source is

\[
\overline{\operatorname{span}\{8F_0\eta_{n,t}:n\ge2\}}^{\mathcal B}
=\overline{F_0\mathcal B}^{\mathcal B}=I,
\qquad F_0\mathcal B\subsetneq I.
\tag{RSR6.1}
\]

Its analytic division input retains the full genus-one product of
\(F_0\), every zero order, the leading exponential and all exponential
factors. Off disks of radius \(R^{-2}\) about zeros of modulus at most
\(8R\), it proves the reciprocal bound

\[
|F_0(s)|^{-1}\le\exp(CR^{3/2}\log R)
\quad(R/2\le|s|\le3R).
\tag{RSR6.2}
\]

The total radii of the zero disks are \(O(R^{-1/2})\). For \(F\in I\),
\(F/F_0\) is entire with all orders removable; maximum modulus on each
disk cluster extends the bound through its interior in a strip enlarged
by one. Thus \(e^{\varepsilon s^2}F/F_0\in\mathcal B\) for every
\(\varepsilon>0\). Its product with \(F_0\) tends to \(F\) by (RSR2.8).
This is the full division mechanism from NCI2 and its independent NER11
verification; it proves a closure, not algebraic equality. The latter
fails already for \(F_0\in I\), since \(F_0=F_0H\) with
\(H\in\mathcal B\) would force \(H=1\), and the constant one is not in
\(\mathcal B\). The discrepancy identification in (RSR6.1) then follows
from (RSR2.9) and continuous multiplication by the full \(8F_0\).
Neither the factor eight nor any zero order is dropped from this
multiplication map.

If \(I\subset P\mathcal B\), the identity on representatives induces

\[
\pi_P:Q\longrightarrow\mathcal B/(P\mathcal B),
\qquad[F]_I\longmapsto[F]_{P\mathcal B},\qquad
\ker\pi_P=P\mathcal B/I.
\tag{RSR6.3}
\]

It is continuous and onto. It is also an open quotient map: for open
\(V\subset Q\), its inverse image under \(\mathcal B\to Q\) is open, and
its image under the open source quotient
\(\mathcal B\to\mathcal B/(P\mathcal B)\) is exactly \(\pi_P(V)\). Thus
the topology asserted in RSS5 is the actual quotient topology. Under
(RSR3.6), its transpose is the continuous inclusion of the selected
finite derivative functionals into \(I^\perp\), retaining (RSR4.4).
Finite-dimensionality makes their inherited Hausdorff locally convex
topology the same as their finite jet topology. No unproved
bounded-lifting claim for an arbitrary infinite-dimensional strong dual
is needed.

\subsection{RSR7. Stronger exact classification of the sections
admitting that
quotient}\label{rsr7.-stronger-exact-classification-of-the-sections-admitting-that-quotient}

RSS5's finite subcollections are exactly all polynomial sections for
which the identity-induced map (RSR6.3) exists. More precisely, for
\(P(0)=1\), the following are equivalent:

\begin{enumerate}
\def\labelenumi{\arabic{enumi}.}
\tightlist
\item
  The rule \([F]_I\mapsto[F]_{P\mathcal B}\) is well-defined.
\item
  \(I\subset P\mathcal B\).
\item
  Every root \(a\) of \(P\) is an actual nontrivial zero of original
  \(\zeta\), and \(d_a\le m_a\).
\end{enumerate}

For the first equivalence, two representatives differ by an arbitrary
element of \(I\), so representative independence is exactly the
containment. For necessity of the third condition, \(F_0\in I\) implies
\(F_0\in P\mathcal B\). Thus \(P\) divides \(F_0\) as an entire
function, and the full zero divisor of (RSR0.3) imposes precisely the
stated points and orders. For sufficiency, any element of \(I\) has all
the root orders required by \(P\); (RSR3.3) therefore puts it in
\(P\mathcal B\). This proves all three implications without any presumed
zero location.

Accordingly, choosing arbitrary roots in a residue section always
defines a legitimate section of \(E\), but its finite jet quotient
receives the original divisor through the identity map exactly under the
proved condition above. This is an exact morphism criterion, not a
statement that other sections are unrelated. For every other section the
shear (RSR5.3--5.9) still compares the full original source and every
original functional.

\subsection{RSR8. Review result and the next calculation actually
performed}\label{rsr8.-review-result-and-the-next-calculation-actually-performed}

RSS0--6 passes the mathematical review in its corrected wording. The
source and its topology, Jensen growth, positive derivative sign,
half-Mellin constants, source integral, parameter continuation,
polynomial division, full jet action, quotient topology, and both shear
signs have been verified above. No RH, zero simplicity, finite actual
divisor or unstated Hardy-domain equality was assumed.

The section choice alone does not prove original geometric purity. The
immediate next question was: what advances the comparison using this
result and the established original ideal? RSR7 answered it by proving
the necessary-and-sufficient criterion for a polynomial section's
quotient to receive the actual original divisor through the identity on
the source. This continues the source-level calculation without changing
the base, assuming a missing vanishing, or replacing the complete
original zeta divisor.

The remaining original target is unchanged: the retained source, its
actual specialization map and the geometric lifting sequence must be
connected by a proved weight argument. Neither the section dependence
proved here nor the unchanged full Hardy receiver supplies that
vanishing by itself. Their strongest established connection is the exact
shear preserving all original functionals, together with the fully
characterized quotient maps (RSR6.3--RSR7).
